% ============================================================================
% CAPÍTULO 1: SISTEMAS BIOLÓGICOS
% Bioinformática para Biologia de Sistemas
% ============================================================================

% ============================================================================
% TEMPLATE BASE PARA APRESENTAÇÕES EM BEAMER
% Bioinformática para Biologia de Sistemas
% ============================================================================

\documentclass[12pt, aspectratio=169]{beamer}

% ============================================================================
% PACOTES ESSENCIAIS
% ============================================================================
\usepackage[utf8]{inputenc}
\usepackage[T1]{fontenc}
\usepackage[brazilian]{babel}
\usepackage{amsmath, amssymb, amsthm}
\usepackage{graphicx}
\usepackage{xcolor}
\usepackage{tikz}
\usetikzlibrary{arrows.meta, shapes, positioning, calc, decorations.pathreplacing, decorations.shapes, decorations.pathmorphing}
\usepackage{listings}
\usepackage{booktabs}
\usepackage{multirow}
\usepackage{hyperref}
\usepackage{chemfig}
\usepackage{chemarr}

% ============================================================================
% CONFIGURAÇÃO DO TEMA
% ============================================================================
\usetheme{Madrid}
\usecolortheme{default}

% Cores personalizadas
\definecolor{azulescuro}{RGB}{0, 51, 102}
\definecolor{azulclaro}{RGB}{51, 153, 255}
\definecolor{verdeacido}{RGB}{102, 204, 0}
\definecolor{laranjacel}{RGB}{255, 153, 51}
\definecolor{roxodna}{RGB}{153, 51, 255}

\setbeamercolor{structure}{fg=azulescuro}
\setbeamercolor{palette primary}{bg=azulescuro,fg=white}
\setbeamercolor{palette secondary}{bg=azulclaro,fg=white}
\setbeamercolor{palette tertiary}{bg=azulescuro,fg=white}
\setbeamercolor{palette quaternary}{bg=azulclaro,fg=white}
\setbeamercolor{block title}{bg=azulescuro,fg=white}
\setbeamercolor{block body}{bg=azulclaro!10,fg=black}
\setbeamercolor{block title alerted}{bg=red!80,fg=white}
\setbeamercolor{block body alerted}{bg=red!10,fg=black}
\setbeamercolor{block title example}{bg=verdeacido!80,fg=white}
\setbeamercolor{block body example}{bg=verdeacido!10,fg=black}

% ============================================================================
% CONFIGURAÇÃO DE CÓDIGO
% ============================================================================
\lstset{
    language=Python,
    basicstyle=\ttfamily\small,
    keywordstyle=\color{azulescuro}\bfseries,
    commentstyle=\color{verdeacido},
    stringstyle=\color{laranjacel},
    numbers=left,
    numberstyle=\tiny\color{gray},
    stepnumber=1,
    numbersep=5pt,
    backgroundcolor=\color{white},
    showspaces=false,
    showstringspaces=false,
    showtabs=false,
    frame=single,
    rulecolor=\color{black},
    tabsize=2,
    captionpos=b,
    breaklines=true,
    breakatwhitespace=false,
    escapeinside={\%*}{*)},
    xleftmargin=10pt,
    xrightmargin=10pt
}

% Definir estilo para R
\lstdefinestyle{Rstyle}{
    language=R,
    basicstyle=\ttfamily\small,
    keywordstyle=\color{azulescuro}\bfseries,
    commentstyle=\color{verdeacido},
    stringstyle=\color{laranjacel}
}

% Definir estilo para MATLAB
\lstdefinestyle{MATLABstyle}{
    language=Matlab,
    basicstyle=\ttfamily\small,
    keywordstyle=\color{azulescuro}\bfseries,
    commentstyle=\color{verdeacido},
    stringstyle=\color{laranjacel}
}

% ============================================================================
% COMANDOS PERSONALIZADOS
% ============================================================================

% Comando para equações destacadas
\newcommand{\destaque}[1]{%
    \begin{alertblock}{}
        \begin{center}
            #1
        \end{center}
    \end{alertblock}
}

% Comando para definições
\newcommand{\definicao}[2]{%
    \begin{block}{Definição: #1}
        #2
    \end{block}
}

% Comando para exemplos
\newcommand{\exemplo}[2]{%
    \begin{exampleblock}{Exemplo: #1}
        #2
    \end{exampleblock}
}

% Comando para observações importantes
\newcommand{\importante}[1]{%
    \begin{alertblock}{Importante!}
        #1
    \end{alertblock}
}

% Comandos matemáticos úteis
\newcommand{\R}{\mathbb{R}}
\newcommand{\N}{\mathbb{N}}
\newcommand{\Z}{\mathbb{Z}}
\newcommand{\dif}{\mathrm{d}}
\newcommand{\parcial}[2]{\frac{\partial #1}{\partial #2}}
\newcommand{\derivada}[2]{\frac{\dif #1}{\dif #2}}
\newcommand{\vect}[1]{\boldsymbol{#1}}

% Comandos biológicos
\newcommand{\gene}[1]{\textit{#1}}
\newcommand{\proteina}[1]{\textsc{#1}}
\newcommand{\especie}[1]{\textit{#1}}

% ============================================================================
% CONFIGURAÇÃO DE RODAPÉ
% ============================================================================
\setbeamertemplate{footline}{
  \leavevmode%
  \hbox{%
  \begin{beamercolorbox}[wd=.33\paperwidth,ht=2.25ex,dp=1ex,center]{author in head/foot}%
    \usebeamerfont{author in head/foot}\insertshortauthor
  \end{beamercolorbox}%
  \begin{beamercolorbox}[wd=.34\paperwidth,ht=2.25ex,dp=1ex,center]{title in head/foot}%
    \usebeamerfont{title in head/foot}\insertshorttitle
  \end{beamercolorbox}%
  \begin{beamercolorbox}[wd=.33\paperwidth,ht=2.25ex,dp=1ex,right]{date in head/foot}%
    \usebeamerfont{date in head/foot}\insertshortdate{}\hspace*{2em}
    \insertframenumber{} / \inserttotalframenumber\hspace*{2ex}
  \end{beamercolorbox}}%
  \vskip0pt%
}

% ============================================================================
% CONFIGURAÇÃO DE NAVEGAÇÃO
% ============================================================================
\setbeamertemplate{navigation symbols}{}

% ============================================================================
% CONFIGURAÇÃO DE ITEMIZE/ENUMERATE
% ============================================================================
\setbeamertemplate{itemize items}[circle]
\setbeamertemplate{enumerate items}[default]

% ============================================================================
% FIM DO TEMPLATE
% ============================================================================


% ============================================================================
% INFORMAÇÕES DO DOCUMENTO
% ============================================================================
\title[Cap. 1: Sistemas Biológicos]{Capítulo 1: Sistemas Biológicos}
\subtitle{Bioinformática para Biologia de Sistemas}
\author{Seu Nome}
\institute{Sua Instituição}
\date{\today}

\begin{document}

% ============================================================================
% SLIDE DE TÍTULO
% ============================================================================
\begin{frame}
\titlepage
\end{frame}

% ============================================================================
% SUMÁRIO
% ============================================================================
\begin{frame}{Sumário}
\tableofcontents
\end{frame}

% ============================================================================
% SEÇÃO 1: INTRODUÇÃO
% ============================================================================
\section{Introdução}

\begin{frame}{Visão Geral do Capítulo}
\begin{block}{Objetivos de Aprendizagem}
\begin{itemize}
    \item Compreender os componentes fundamentais dos sistemas biológicos
    \item Entender a estrutura e função do DNA, RNA e proteínas
    \item Explorar processos metabólicos básicos
    \item Conhecer a organização celular
\end{itemize}
\end{block}

\vspace{0.5cm}

\begin{exampleblock}{Por que Estudar Sistemas Biológicos?}
A compreensão dos sistemas biológicos é fundamental para:
\begin{itemize}
    \item Modelagem matemática de processos biológicos
    \item Desenvolvimento de terapias e medicamentos
    \item Engenharia de sistemas biológicos
\end{itemize}
\end{exampleblock}
\end{frame}

\begin{frame}{O que são Sistemas Biológicos?}
\definicao{Sistema Biológico}{
Um sistema biológico é uma rede complexa de componentes biologicamente relevantes que interagem entre si para realizar funções específicas.
}

\vspace{0.5cm}

\begin{columns}
\column{0.5\textwidth}
\textbf{Características:}
\begin{itemize}
    \item Complexidade hierárquica
    \item Emergência de propriedades
    \item Autorregulação
    \item Adaptabilidade
\end{itemize}

\column{0.5\textwidth}
\textbf{Níveis de Organização:}
\begin{itemize}
    \item Molecular (DNA, proteínas)
    \item Celular (organelas, células)
    \item Tecidual
    \item Organismo completo
\end{itemize}
\end{columns}
\end{frame}

% ============================================================================
% SEÇÃO 2: COMPONENTES BIOLÓGICOS
% ============================================================================
\section{Componentes Biológicos}

\begin{frame}{1.1 Componentes Biológicos}
\begin{block}{Os Quatro Pilares da Vida}
\begin{enumerate}
    \item \textbf{Ácidos Nucleicos} (DNA, RNA)
    \item \textbf{Proteínas} (enzimas, estruturais, sinalizadoras)
    \item \textbf{Lipídios} (membranas, armazenamento)
    \item \textbf{Carboidratos} (energia, estrutura)
\end{enumerate}
\end{block}

\vspace{0.3cm}

\begin{columns}
\column{0.5\textwidth}
\begin{exampleblock}{Moléculas Pequenas}
\begin{itemize}
    \item Nucleotídeos
    \item Aminoácidos
    \item Monossacarídeos
    \item Ácidos graxos
\end{itemize}
\end{exampleblock}

\column{0.5\textwidth}
\begin{exampleblock}{Macromoléculas}
\begin{itemize}
    \item DNA/RNA
    \item Proteínas
    \item Polissacarídeos
    \item Lipídios complexos
\end{itemize}
\end{exampleblock}
\end{columns}
\end{frame}

\begin{frame}{Dogma Central da Biologia Molecular}
\begin{center}
\begin{tikzpicture}[node distance=3cm, auto,
    box/.style={rectangle, draw, fill=azulclaro!20, text width=2cm, text centered, rounded corners, minimum height=1cm}]

    \node[box, fill=roxodna!20] (dna) {DNA};
    \node[box, fill=laranjacel!20, right of=dna] (rna) {RNA};
    \node[box, fill=verdeacido!20, right of=rna] (prot) {Proteína};

    \draw[->, thick, azulescuro] (dna) -- node[above] {Transcrição} (rna);
    \draw[->, thick, azulescuro] (rna) -- node[above] {Tradução} (prot);
    \draw[->, thick, azulescuro] (dna) to [out=135,in=45,looseness=8] node[above] {Replicação} (dna);

\end{tikzpicture}
\end{center}

\importante{
O fluxo de informação genética é unidirecional: DNA $\rightarrow$ RNA $\rightarrow$ Proteína
}
\end{frame}

% ============================================================================
% SEÇÃO 3: ESTRUTURA DO DNA
% ============================================================================
\section{Estrutura do DNA}

\begin{frame}{1.2 Estrutura do DNA}
\begin{columns}
\column{0.5\textwidth}
\textbf{Características Estruturais:}
\begin{itemize}
    \item Dupla hélice
    \item Fitas antiparalelas
    \item Pareamento de bases
    \item Sulco maior e menor
\end{itemize}

\vspace{0.3cm}

\destaque{
$\text{A} \equiv \text{T} \quad \text{(2 ligações H)}$\\
$\text{G} \equiv \text{C} \quad \text{(3 ligações H)}$
}

\column{0.5\textwidth}
\begin{center}
\begin{tikzpicture}[scale=0.8]
    % Backbone esquerdo
    \draw[thick, azulescuro] (0,0) -- (0,4);
    % Backbone direito
    \draw[thick, azulescuro] (3,0) -- (3,4);

    % Pares de bases
    \draw[thick, roxodna] (0,0.5) -- node[above, font=\tiny] {A-T} (3,0.5);
    \draw[thick, verdeacido] (0,1.5) -- node[above, font=\tiny] {G-C} (3,1.5);
    \draw[thick, roxodna] (0,2.5) -- node[above, font=\tiny] {T-A} (3,2.5);
    \draw[thick, verdeacido] (0,3.5) -- node[above, font=\tiny] {C-G} (3,3.5);

    % Setas de direção
    \draw[->, thick] (-0.3,0) -- (-0.3,1) node[midway, left, font=\tiny] {5'};
    \draw[->, thick] (3.3,4) -- (3.3,3) node[midway, right, font=\tiny] {3'};
\end{tikzpicture}
\end{center}
\end{columns}
\end{frame}

\begin{frame}{Composição Química do DNA}
\definicao{Nucleotídeo}{
Unidade básica do DNA composta por três componentes:
\begin{enumerate}
    \item Base nitrogenada (A, T, G, C)
    \item Açúcar (desoxirribose)
    \item Grupo fosfato
\end{enumerate}
}

\vspace{0.3cm}

\begin{block}{Estrutura Química}
\begin{equation*}
\text{Nucleotídeo} = \underbrace{\text{Base}}_{\text{Purinas/Pirimidinas}} + \underbrace{\text{Desoxirribose}}_{\text{Açúcar}} + \underbrace{\text{PO}_4^{3-}}_{\text{Fosfato}}
\end{equation*}
\end{block}

\begin{exampleblock}{Purinas vs Pirimidinas}
\begin{itemize}
    \item \textbf{Purinas:} Adenina (A), Guanina (G) --- dois anéis
    \item \textbf{Pirimidinas:} Timina (T), Citosina (C) --- um anel
\end{itemize}
\end{exampleblock}
\end{frame}

\begin{frame}{Propriedades Físico-Químicas do DNA}
\begin{block}{Parâmetros Estruturais da Dupla Hélice}
\begin{itemize}
    \item \textbf{Diâmetro:} $\sim$2 nm
    \item \textbf{Comprimento por volta:} 3.4 nm
    \item \textbf{Pares de bases por volta:} 10 pb
    \item \textbf{Ângulo de rotação:} 36° por pb
\end{itemize}
\end{block}

\vspace{0.3cm}

\begin{columns}
\column{0.5\textwidth}
\begin{exampleblock}{Formas do DNA}
\begin{itemize}
    \item \textbf{B-DNA:} forma fisiológica
    \item \textbf{A-DNA:} desidratado
    \item \textbf{Z-DNA:} hélice esquerda
\end{itemize}
\end{exampleblock}

\column{0.5\textwidth}
\begin{alertblock}{Desnaturação}
\begin{itemize}
    \item Temperatura
    \item pH extremos
    \item Agentes químicos
    \item $T_m$ = temperatura de fusão
\end{itemize}
\end{alertblock}
\end{columns}
\end{frame}

% ============================================================================
% SEÇÃO 4: ESTRUTURA DO RNA
% ============================================================================
\section{Estrutura do RNA}

\begin{frame}{1.3 Estrutura do RNA}
\begin{block}{Diferenças DNA vs RNA}
\begin{center}
\begin{tabular}{lcc}
\toprule
\textbf{Característica} & \textbf{DNA} & \textbf{RNA} \\
\midrule
Açúcar & Desoxirribose & Ribose \\
Bases & A, T, G, C & A, U, G, C \\
Estrutura & Dupla fita & Fita simples* \\
Estabilidade & Alta & Menor \\
Tamanho & Longo & Variável \\
\bottomrule
\end{tabular}
\end{center}
\end{block}

\vspace{0.2cm}

\begin{exampleblock}{* Estrutura do RNA}
Embora seja fita simples, o RNA pode formar estruturas secundárias e terciárias complexas por pareamento intramolecular.
\end{exampleblock}
\end{frame}

\begin{frame}{Tipos de RNA}
\begin{columns}
\column{0.5\textwidth}
\begin{block}{RNA Mensageiro (mRNA)}
\begin{itemize}
    \item Carrega informação genética
    \item Traduzido em proteínas
    \item 5\% do RNA celular
    \item Possui códons
\end{itemize}
\end{block}

\begin{block}{RNA Transportador (tRNA)}
\begin{itemize}
    \item Transporta aminoácidos
    \item Estrutura em trevo
    \item Anticódon
    \item 15\% do RNA celular
\end{itemize}
\end{block}

\column{0.5\textwidth}
\begin{block}{RNA Ribossomal (rRNA)}
\begin{itemize}
    \item Componente dos ribossomos
    \item Atividade catalítica
    \item 80\% do RNA celular
    \item Altamente conservado
\end{itemize}
\end{block}

\begin{block}{RNAs Regulatórios}
\begin{itemize}
    \item microRNA (miRNA)
    \item RNA interferente (siRNA)
    \item RNA longo não-codificante
    \item RNA ribossômico pequeno
\end{itemize}
\end{block}
\end{columns}
\end{frame}

\begin{frame}{Estrutura Secundária do RNA}
\begin{columns}
\column{0.5\textwidth}
\textbf{Elementos Estruturais:}
\begin{itemize}
    \item Hairpin loops (grampos)
    \item Bulges (protuberâncias)
    \item Internal loops (alças internas)
    \item Pseudoknots (pseudonós)
\end{itemize}

\vspace{0.3cm}

\definicao{Estrutura Secundária}{
Arranjo espacial resultante do pareamento de bases complementares dentro da mesma molécula de RNA.
}

\column{0.5\textwidth}
\begin{center}
\begin{tikzpicture}[scale=0.7]
    % Stem-loop structure
    \draw[thick, azulescuro] (0,0) -- (0,2);
    \draw[thick, azulescuro] (1,0) -- (1,2);
    \draw[thick, azulescuro] (0,2) to [out=0, in=180] (1,2);

    % Base pairs
    \foreach \y in {0.3, 0.7, 1.1, 1.5}
        \draw[roxodna] (0,\y) -- (1,\y);

    % Labels
    \node[left, font=\tiny] at (0,1) {5'};
    \node[right, font=\tiny] at (1,1) {3'};
    \node[above, font=\small] at (0.5,2.5) {Hairpin};

    % Bulge example
    \begin{scope}[xshift=3cm]
        \draw[thick, verdeacido] (0,0) -- (0,0.5);
        \draw[thick, verdeacido] (1,0) -- (1,1);
        \draw[thick, verdeacido] (0,0.5) -- (-0.3,1);
        \draw[thick, verdeacido] (-0.3,1) -- (0,1.5);
        \draw[thick, verdeacido] (1,1) -- (1,1.5);

        \node[above, font=\small] at (0.5,1.8) {Bulge};
    \end{scope}
\end{tikzpicture}
\end{center}
\end{columns}
\end{frame}

\begin{frame}[fragile]{Predição de Estrutura Secundária do RNA}
\begin{exampleblock}{Algoritmo de Nussinov}
Algoritmo de programação dinâmica para predizer estrutura secundária de RNA:
\end{exampleblock}

\begin{lstlisting}[language=Python, caption=Pseudocódigo para Nussinov, basicstyle=\ttfamily\footnotesize]
def nussinov(sequence):
    n = len(sequence)
    dp = [[0] * n for _ in range(n)]

    for length in range(2, n + 1):
        for i in range(n - length + 1):
            j = i + length - 1
            # Caso 1: i e j não pareiam
            dp[i][j] = dp[i+1][j]

            # Caso 2: i e j pareiam
            for k in range(i, j):
                if can_pair(sequence[i], sequence[k]):
                    score = dp[i+1][k-1] + 1 + dp[k+1][j]
                    dp[i][j] = max(dp[i][j], score)

    return dp[0][n-1]
\end{lstlisting}
\end{frame}

% ============================================================================
% SEÇÃO 5: ESTRUTURA DE PROTEÍNAS
% ============================================================================
\section{Estrutura de Proteínas}

\begin{frame}{1.4 Estrutura de Proteínas}
\begin{block}{Níveis de Organização Proteica}
\begin{enumerate}
    \item \textbf{Estrutura Primária:} sequência de aminoácidos
    \item \textbf{Estrutura Secundária:} $\alpha$-hélices, $\beta$-folhas
    \item \textbf{Estrutura Terciária:} dobramento 3D completo
    \item \textbf{Estrutura Quaternária:} complexos de múltiplas cadeias
\end{enumerate}
\end{block}

\vspace{0.3cm}

\begin{center}
\begin{tikzpicture}[scale=0.8]
    \node[draw, rectangle, fill=azulclaro!20, text width=2cm, align=center] at (0,0) {Primária\\(1D)};
    \node[draw, rectangle, fill=azulclaro!40, text width=2cm, align=center] at (3,0) {Secundária\\(2D)};
    \node[draw, rectangle, fill=azulclaro!60, text width=2cm, align=center] at (6,0) {Terciária\\(3D)};
    \node[draw, rectangle, fill=azulclaro!80, text width=2cm, align=center] at (9,0) {Quaternária\\(Complexo)};

    \draw[->, thick] (1,0) -- (2,0);
    \draw[->, thick] (4,0) -- (5,0);
    \draw[->, thick] (7,0) -- (8,0);
\end{tikzpicture}
\end{center}
\end{frame}

\begin{frame}{Aminoácidos: Blocos Construtores das Proteínas}
\begin{columns}
\column{0.5\textwidth}
\begin{block}{Estrutura Geral}
\begin{center}
\chemfig{H_2N-C(-[2]H)(-[6]R)-C(=[1]O)-[7]OH}
\end{center}
\begin{itemize}
    \item Grupo amino (NH$_2$)
    \item Carbono $\alpha$
    \item Cadeia lateral (R)
    \item Grupo carboxila (COOH)
\end{itemize}
\end{block}

\column{0.5\textwidth}
\begin{exampleblock}{Classificação dos Aminoácidos}
\begin{itemize}
    \item \textbf{Apolar:} Ala, Val, Leu, Ile, Met, Pro, Phe, Trp
    \item \textbf{Polar:} Ser, Thr, Cys, Tyr, Asn, Gln
    \item \textbf{Ácido:} Asp, Glu
    \item \textbf{Básico:} Lys, Arg, His
\end{itemize}
\end{exampleblock}
\end{columns}

\vspace{0.3cm}

\importante{
Total de 20 aminoácidos padrão codificados geneticamente
}
\end{frame}

\begin{frame}{Ligação Peptídica}
\begin{columns}
\column{0.6\textwidth}
\begin{block}{Formação da Ligação Peptídica}
Reação de condensação entre dois aminoácidos:
\begin{equation*}
\text{AA}_1\text{-COOH} + \text{H-NH-AA}_2 \rightarrow \text{AA}_1\text{-CO-NH-AA}_2 + \text{H}_2\text{O}
\end{equation*}
\end{block}

\textbf{Características:}
\begin{itemize}
    \item Ligação covalente
    \item Caráter parcial de dupla ligação
    \item Planar e rígida
    \item Trans configuração preferida
\end{itemize}

\column{0.4\textwidth}
\begin{alertblock}{Ângulos de Torção}
\begin{itemize}
    \item $\phi$ (phi): N-C$_\alpha$
    \item $\psi$ (psi): C$_\alpha$-C
    \item $\omega$ (omega): ligação peptídica ($\sim$180°)
\end{itemize}
\end{alertblock}

\begin{exampleblock}{Diagrama de Ramachandran}
Representa regiões permitidas dos ângulos $\phi$ e $\psi$
\end{exampleblock}
\end{columns}
\end{frame}

\begin{frame}{Estrutura Secundária}
\begin{columns}
\column{0.5\textwidth}
\begin{block}{$\alpha$-Hélice}
\begin{itemize}
    \item Hélice destrogira
    \item 3.6 resíduos/volta
    \item Ligações H: C=O(i) $\leftrightarrow$ N-H(i+4)
    \item Compacta e estável
\end{itemize}

\textbf{Parâmetros:}
\begin{itemize}
    \item $\phi = -60°$
    \item $\psi = -45°$
    \item Passo: 5.4 Å
\end{itemize}
\end{block}

\column{0.5\textwidth}
\begin{block}{Folha $\beta$}
\begin{itemize}
    \item Estrutura estendida
    \item Fitas paralelas ou antiparalelas
    \item Ligações H entre fitas
    \item Pregueada
\end{itemize}

\textbf{Tipos:}
\begin{itemize}
    \item Paralela: mesma direção
    \item Antiparalela: direções opostas
    \item Mista
\end{itemize}
\end{block}
\end{columns}

\vspace{0.3cm}

\begin{exampleblock}{Outras Estruturas}
Turns (voltas), loops (alças), coils (espirais irregulares)
\end{exampleblock}
\end{frame}

\begin{frame}{Estrutura Terciária e Quaternária}
\begin{block}{Estrutura Terciária}
Arranjo tridimensional completo de uma cadeia polipeptídica
\begin{itemize}
    \item \textbf{Interações:} ligações de hidrogênio, iônicas, hidrofóbicas, pontes dissulfeto
    \item \textbf{Domínios:} unidades estruturais e funcionais
    \item \textbf{Motifs:} padrões estruturais recorrentes
\end{itemize}
\end{block}

\vspace{0.3cm}

\begin{block}{Estrutura Quaternária}
Associação de múltiplas cadeias polipeptídicas (subunidades)
\end{block}

\exemplo{Hemoglobina}{
Proteína tetramérica: 2 subunidades $\alpha$ + 2 subunidades $\beta$\\
Função: transporte de O$_2$ no sangue
}
\end{frame}

\begin{frame}[fragile]{Bancos de Dados de Proteínas}
\begin{block}{Protein Data Bank (PDB)}
Principal repositório de estruturas 3D de proteínas
\begin{itemize}
    \item URL: \url{https://www.rcsb.org}
    \item Mais de 200.000 estruturas
    \item Formato PDB, mmCIF
    \item Técnicas: Cristalografia de raios-X, NMR, Cryo-EM
\end{itemize}
\end{block}

\begin{exampleblock}{Acessando PDB com BioPython}
\begin{lstlisting}[language=Python]
from Bio.PDB import PDBParser

parser = PDBParser()
structure = parser.get_structure("1A3N", "1a3n.pdb")

for model in structure:
    for chain in model:
        for residue in chain:
            print(f"Residue: {residue.resname}")
\end{lstlisting}
\end{exampleblock}
\end{frame}

% ============================================================================
% SEÇÃO 6: PROCESSOS METABÓLICOS
% ============================================================================
\section{Processos Metabólicos}

\begin{frame}{1.5 Processos Metabólicos}
\definicao{Metabolismo}{
Conjunto de reações químicas que ocorrem nas células para manter a vida, dividido em:
\begin{itemize}
    \item \textbf{Catabolismo:} degradação de moléculas (libera energia)
    \item \textbf{Anabolismo:} síntese de moléculas (consome energia)
\end{itemize}
}

\vspace{0.3cm}

\begin{center}
\begin{tikzpicture}[node distance=2cm]
    \node[draw, circle, fill=laranjacel!30, text width=2cm, align=center] (cat) {Catabolismo};
    \node[draw, circle, fill=verdeacido!30, text width=2cm, align=center, right=3cm of cat] (ana) {Anabolismo};

    \draw[->, thick, bend left=30] (cat) to node[above] {ATP, NADH} (ana);
    \draw[->, thick, bend left=30] (ana) to node[below] {ADP, NAD$^+$} (cat);

    \node[above=0.3cm of cat, font=\small] {Moléculas\\complexas};
    \node[below=0.3cm of cat, font=\small] {Moléculas\\simples};
    \node[above=0.3cm of ana, font=\small] {Moléculas\\simples};
    \node[below=0.3cm of ana, font=\small] {Moléculas\\complexas};
\end{tikzpicture}
\end{center}
\end{frame}

\begin{frame}{Vias Metabólicas Principais}
\begin{columns}
\column{0.5\textwidth}
\begin{block}{Metabolismo de Carboidratos}
\begin{itemize}
    \item \textbf{Glicólise}
    \item Ciclo de Krebs (TCA)
    \item Fosforilação oxidativa
    \item Gliconeogênese
    \item Via das pentoses
\end{itemize}
\end{block}

\begin{block}{Metabolismo de Lipídios}
\begin{itemize}
    \item $\beta$-oxidação
    \item Síntese de ácidos graxos
    \item Síntese de colesterol
\end{itemize}
\end{block}

\column{0.5\textwidth}
\begin{block}{Metabolismo de Aminoácidos}
\begin{itemize}
    \item Transaminação
    \item Desaminação
    \item Ciclo da ureia
    \item Síntese de aa
\end{itemize}
\end{block}

\begin{block}{Metabolismo de Nucleotídeos}
\begin{itemize}
    \item Síntese de purinas
    \item Síntese de pirimidinas
    \item Vias de salvação
\end{itemize}
\end{block}
\end{columns}
\end{frame}

\begin{frame}{Glicólise: Via Fundamental}
\begin{block}{Visão Geral}
\begin{equation*}
\text{Glicose} + 2\text{NAD}^+ + 2\text{ADP} + 2\text{P}_i \rightarrow 2\text{Piruvato} + 2\text{NADH} + 2\text{ATP} + 2\text{H}_2\text{O}
\end{equation*}
\end{block}

\begin{columns}
\column{0.5\textwidth}
\textbf{Fase Preparatória} (consome ATP):
\begin{enumerate}
    \item Glicose $\rightarrow$ G6P
    \item G6P $\rightarrow$ F6P
    \item F6P $\rightarrow$ F1,6BP
    \item F1,6BP $\rightarrow$ DHAP + G3P
    \item DHAP $\rightarrow$ G3P
\end{enumerate}

\column{0.5\textwidth}
\textbf{Fase de Pagamento} (produz ATP):
\begin{enumerate}
    \setcounter{enumi}{5}
    \item G3P $\rightarrow$ 1,3BPG
    \item 1,3BPG $\rightarrow$ 3PG
    \item 3PG $\rightarrow$ 2PG
    \item 2PG $\rightarrow$ PEP
    \item PEP $\rightarrow$ Piruvato
\end{enumerate}
\end{columns}

\vspace{0.3cm}

\importante{
Balanço líquido: 2 ATP, 2 NADH produzidos por molécula de glicose
}
\end{frame}

\begin{frame}{Enzimas: Catalisadores Biológicos}
\definicao{Enzima}{
Proteína que catalisa (acelera) reações bioquímicas específicas sem ser consumida no processo.
}

\vspace{0.3cm}

\begin{block}{Cinética de Michaelis-Menten}
\begin{equation*}
v = \frac{V_{max}[S]}{K_M + [S]}
\end{equation*}

Onde:
\begin{itemize}
    \item $v$ = velocidade da reação
    \item $V_{max}$ = velocidade máxima
    \item $[S]$ = concentração do substrato
    \item $K_M$ = constante de Michaelis (afinidade)
\end{itemize}
\end{block}

\begin{exampleblock}{Interpretação}
Quando $[S] = K_M$, então $v = V_{max}/2$
\end{exampleblock}
\end{frame}

\begin{frame}[fragile]{Modelando Cinética Enzimática}
\begin{lstlisting}[language=Python, caption=Simulação de Michaelis-Menten, basicstyle=\ttfamily\footnotesize]
import numpy as np
import matplotlib.pyplot as plt

def michaelis_menten(S, Vmax, Km):
    """Calcula velocidade pela equacao de MM"""
    return (Vmax * S) / (Km + S)

# Parametros
Vmax = 100  # umol/min
Km = 5      # mM

# Concentracoes de substrato
S = np.linspace(0, 50, 100)
v = michaelis_menten(S, Vmax, Km)

# Plotar
plt.plot(S, v, 'b-', linewidth=2)
plt.axhline(y=Vmax, color='r', linestyle='--',
            label=f'Vmax = {Vmax}')
plt.axhline(y=Vmax/2, color='g', linestyle='--')
plt.axvline(x=Km, color='g', linestyle='--',
            label=f'Km = {Km}')
plt.xlabel('[S] (mM)')
plt.ylabel('v (umol/min)')
plt.legend()
\end{lstlisting}
\end{frame}

% ============================================================================
% SEÇÃO 7: ORGANIZAÇÃO CELULAR
% ============================================================================
\section{Organização Celular}

\begin{frame}{1.6 Organização Celular}
\begin{columns}
\column{0.5\textwidth}
\begin{block}{Célula Procarionte}
\begin{itemize}
    \item Sem núcleo definido
    \item DNA circular
    \item Sem organelas membranosas
    \item Ribossomos 70S
    \item Parede celular
\end{itemize}
\textbf{Exemplos:} Bactérias, Archaea
\end{block}

\column{0.5\textwidth}
\begin{block}{Célula Eucarionte}
\begin{itemize}
    \item Núcleo definido
    \item DNA linear + histonas
    \item Organelas membranosas
    \item Ribossomos 80S
    \item Citoesqueleto
\end{itemize}
\textbf{Exemplos:} Animais, Plantas, Fungos
\end{block}
\end{columns}

\vspace{0.5cm}

\begin{center}
\begin{tabular}{lcc}
\toprule
\textbf{Característica} & \textbf{Procarionte} & \textbf{Eucarionte} \\
\midrule
Tamanho típico & 1-10 $\mu$m & 10-100 $\mu$m \\
Complexidade & Simples & Complexa \\
\bottomrule
\end{tabular}
\end{center}
\end{frame}

\begin{frame}{Organelas Celulares}
\begin{columns}
\column{0.5\textwidth}
\begin{block}{Organelas Principais}
\begin{itemize}
    \item \textbf{Núcleo:} DNA, transcrição
    \item \textbf{Mitocôndria:} ATP, respiração
    \item \textbf{RE Rugoso:} síntese proteica
    \item \textbf{RE Liso:} lipídios
    \item \textbf{Golgi:} processamento
    \item \textbf{Lisossomo:} digestão
\end{itemize}
\end{block}

\column{0.5\textwidth}
\begin{block}{Organelas em Plantas}
\begin{itemize}
    \item \textbf{Cloroplasto:} fotossíntese
    \item \textbf{Vacúolo:} armazenamento
    \item \textbf{Parede celular:} suporte
\end{itemize}
\end{block}

\begin{exampleblock}{Teoria Endossimbiótica}
Mitocôndrias e cloroplastos originaram-se de bactérias simbiontes
\end{exampleblock}
\end{columns}
\end{frame}

\begin{frame}{Mitocôndria: Usina de Energia}
\begin{columns}
\column{0.6\textwidth}
\textbf{Estrutura:}
\begin{itemize}
    \item Membrana externa (permeável)
    \item Membrana interna (pregueada - cristas)
    \item Matriz mitocondrial
    \item Espaço intermembranar
    \item DNA mitocondrial próprio
\end{itemize}

\vspace{0.3cm}

\textbf{Funções:}
\begin{itemize}
    \item Ciclo de Krebs (matriz)
    \item Cadeia respiratória (membrana interna)
    \item Síntese de ATP
    \item Metabolismo de lipídios
    \item Apoptose (morte celular programada)
\end{itemize}

\column{0.4\textwidth}
\begin{alertblock}{Produção de ATP}
\begin{equation*}
\begin{split}
&\text{Glicólise: } 2 \text{ ATP}\\
&\text{Krebs: } 2 \text{ ATP}\\
&\text{Fosf. Ox: } 34 \text{ ATP}\\
&\hline
&\textbf{Total: } 38 \text{ ATP}
\end{split}
\end{equation*}
\end{alertblock}

\begin{exampleblock}{Herança Materna}
DNA mitocondrial é herdado apenas da mãe
\end{exampleblock}
\end{columns}
\end{frame}

\begin{frame}{Compartimentalização Celular}
\importante{
A compartimentalização permite:
\begin{itemize}
    \item Concentração de enzimas e substratos
    \item Ambientes químicos distintos (pH, íons)
    \item Processos simultâneos incompatíveis
    \item Regulação espacial do metabolismo
\end{itemize}
}

\vspace{0.3cm}

\begin{exampleblock}{Exemplo: Síntese vs Degradação}
\begin{itemize}
    \item \textbf{Citoplasma:} síntese de ácidos graxos
    \item \textbf{Mitocôndria:} $\beta$-oxidação de ácidos graxos
    \item Ambos podem ocorrer simultaneamente sem interferência
\end{itemize}
\end{exampleblock}
\end{frame}

\begin{frame}{Transporte através de Membranas}
\begin{block}{Tipos de Transporte}
\begin{enumerate}
    \item \textbf{Passivo:} a favor do gradiente (sem ATP)
    \begin{itemize}
        \item Difusão simples
        \item Difusão facilitada (canais, transportadores)
    \end{itemize}

    \item \textbf{Ativo:} contra o gradiente (consome ATP)
    \begin{itemize}
        \item Primário: bombas (Na$^+$/K$^+$-ATPase)
        \item Secundário: cotransporte
    \end{itemize}

    \item \textbf{Em massa:}
    \begin{itemize}
        \item Endocitose (para dentro)
        \item Exocitose (para fora)
    \end{itemize}
\end{enumerate}
\end{block}
\end{frame}

\begin{frame}{Equação de Nernst: Potencial de Membrana}
\begin{block}{Potencial de Equilíbrio}
Para um íon específico:
\begin{equation*}
E_{ion} = \frac{RT}{zF} \ln\frac{[ion]_{ext}}{[ion]_{int}} = \frac{58\text{ mV}}{z} \log_{10}\frac{[ion]_{ext}}{[ion]_{int}}
\end{equation*}

Onde:
\begin{itemize}
    \item $R$ = constante dos gases
    \item $T$ = temperatura (K)
    \item $z$ = valência do íon
    \item $F$ = constante de Faraday
\end{itemize}
\end{block}

\exemplo{Potencial de K$^+$}{
Se $[K^+]_{ext} = 5$ mM e $[K^+]_{int} = 140$ mM:\\
$E_{K^+} = 58 \log_{10}(5/140) \approx -85$ mV
}
\end{frame}

% ============================================================================
% SEÇÃO 8: INTEGRAÇÃO DE SISTEMAS
% ============================================================================
\section{Integração de Sistemas}

\begin{frame}{Visão Integrada dos Sistemas Biológicos}
\begin{center}
\begin{tikzpicture}[node distance=1.5cm, scale=0.9, transform shape]
    \node[draw, circle, fill=roxodna!30] (dna) {DNA};
    \node[draw, circle, fill=laranjacel!30, right=of dna] (rna) {RNA};
    \node[draw, circle, fill=verdeacido!30, right=of rna] (prot) {Proteína};
    \node[draw, circle, fill=azulclaro!30, below=of prot] (met) {Metabólito};
    \node[draw, circle, fill=yellow!30, below=of rna] (cel) {Célula};

    \draw[->, thick] (dna) -- (rna);
    \draw[->, thick] (rna) -- (prot);
    \draw[->, thick] (prot) -- (met);
    \draw[->, thick] (met) -- (cel);
    \draw[->, thick, bend right=45] (prot) to node[left, font=\tiny] {regulação} (dna);
    \draw[->, thick, bend left=20] (met) to node[above right, font=\tiny] {feedback} (prot);
    \draw[->, thick] (cel) -- ($(cel) + (0,-1)$) node[below, font=\small] {Fenótipo};
\end{tikzpicture}
\end{center}

\importante{
Os sistemas biológicos são redes altamente interconectadas onde:
\begin{itemize}
    \item Genes regulam proteínas
    \item Proteínas catalisam metabolismo
    \item Metabólitos regulam genes
    \item Tudo ocorre no contexto celular
\end{itemize}
}
\end{frame}

\begin{frame}{Da Molécula ao Organismo}
\begin{block}{Hierarquia de Organização Biológica}
\begin{center}
\begin{tikzpicture}[scale=0.8]
    \foreach \i/\name in {0/Átomo, 1/Molécula, 2/Macromolécula, 3/Organela, 4/Célula, 5/Tecido, 6/Órgão, 7/Organismo} {
        \node[draw, rectangle, fill=azulclaro!20, minimum width=2.5cm] at (0,-\i*0.8) {\name};
    }

    \foreach \i in {0,...,6} {
        \draw[->, thick] (-1.3,-\i*0.8-0.3) -- (-1.3,-\i*0.8-0.5);
    }

    \node[right, font=\small] at (1.5,-3.2) {Emergência};
    \node[right, font=\small] at (1.5,-3.5) {de};
    \node[right, font=\small] at (1.5,-3.8) {propriedades};
\end{tikzpicture}
\end{center}
\end{block}
\end{frame}

\begin{frame}{Biologia de Sistemas: Abordagem Holística}
\definicao{Biologia de Sistemas}{
Estudo de sistemas biológicos através da análise integrada de múltiplos níveis de organização (genoma, transcriptoma, proteoma, metaboloma).
}

\vspace{0.3cm}

\begin{columns}
\column{0.5\textwidth}
\textbf{Abordagem Reducionista:}
\begin{itemize}
    \item Estudo de componentes isolados
    \item Foco em detalhes
    \item Experimentos controlados
\end{itemize}

\column{0.5\textwidth}
\textbf{Abordagem Sistêmica:}
\begin{itemize}
    \item Estudo de interações
    \item Visão do todo
    \item Modelagem integrativa
\end{itemize}
\end{columns}

\vspace{0.3cm}

\importante{
A biologia de sistemas combina ambas as abordagens!
}
\end{frame}

% ============================================================================
% SEÇÃO 9: EXERCÍCIOS
% ============================================================================
\section{Exercícios}

\begin{frame}{Exercícios - Parte 1}
\begin{block}{Questão 1: Estrutura do DNA}
Calcule o número total de ligações de hidrogênio em uma molécula de DNA de 100 pares de bases, sabendo que ela contém 30\% de adenina.
\end{block}

\begin{block}{Questão 2: Cinética Enzimática}
Uma enzima tem $K_M = 2$ mM e $V_{max} = 50$ $\mu$mol/min. Calcule:
\begin{enumerate}
    \item A velocidade quando $[S] = 2$ mM
    \item A concentração de substrato necessária para atingir 90\% de $V_{max}$
\end{enumerate}
\end{block}

\begin{block}{Questão 3: Código Genético}
Quantas proteínas diferentes de 5 aminoácidos podem ser formadas teoricamente usando os 20 aminoácidos padrão?
\end{block}
\end{frame}

\begin{frame}{Exercícios - Parte 2}
\begin{block}{Questão 4: Estrutura de Proteínas}
Explique como cada nível de estrutura proteica (primária, secundária, terciária, quaternária) contribui para a função da hemoglobina.
\end{block}

\begin{block}{Questão 5: Metabolismo}
Calcule o rendimento total de ATP a partir da oxidação completa de uma molécula de glicose, considerando:
\begin{itemize}
    \item Glicólise: 2 ATP + 2 NADH
    \item Conversão piruvato $\rightarrow$ Acetil-CoA: 2 NADH
    \item Ciclo de Krebs: 2 ATP + 6 NADH + 2 FADH$_2$
    \item Cada NADH rende 2.5 ATP; cada FADH$_2$ rende 1.5 ATP
\end{itemize}
\end{block}
\end{frame}

\begin{frame}{Exercícios - Parte 3}
\begin{block}{Questão 6: Potencial de Membrana}
Calcule o potencial de equilíbrio para Na$^+$ a 37°C, dados:
\begin{itemize}
    \item $[Na^+]_{ext} = 145$ mM
    \item $[Na^+]_{int} = 12$ mM
\end{itemize}
Use a equação de Nernst simplificada: $E = (58/z) \log([X]_{ext}/[X]_{int})$
\end{block}

\begin{alertblock}{Projeto Prático}
Escolha uma via metabólica e:
\begin{enumerate}
    \item Desenhe o diagrama completo com todas as enzimas
    \item Identifique pontos de regulação
    \item Proponha um modelo matemático simples para uma reação chave
    \item Implemente o modelo em Python ou R
\end{enumerate}
\end{alertblock}
\end{frame}

% ============================================================================
% SEÇÃO 10: REFERÊNCIAS
% ============================================================================
\section{Referências}

\begin{frame}{Referências Principais}
\begin{thebibliography}{99}
\bibitem{alberts} Alberts B. et al. (2022). \textit{Molecular Biology of the Cell}, 7th ed. Garland Science.

\bibitem{berg} Berg J.M., Tymoczko J.L., Stryer L. (2019). \textit{Biochemistry}, 9th ed. W.H. Freeman.

\bibitem{nelson} Nelson D.L., Cox M.M. (2021). \textit{Lehninger Principles of Biochemistry}, 8th ed. W.H. Freeman.

\bibitem{lodish} Lodish H. et al. (2021). \textit{Molecular Cell Biology}, 9th ed. W.H. Freeman.

\bibitem{palsson} Palsson B.O., Zengler K. (2010). \textit{The challenges of integrating multi-omic data sets}. Nature Chemical Biology.
\end{thebibliography}
\end{frame}

\begin{frame}{Leitura Complementar}
\begin{block}{Livros Recomendados}
\begin{itemize}
    \item Klipp E. et al. (2016). \textit{Systems Biology: A Textbook}, 2nd ed.
    \item Alon U. (2019). \textit{An Introduction to Systems Biology}, 2nd ed.
    \item Voit E.O. (2017). \textit{A First Course in Systems Biology}, 2nd ed.
\end{itemize}
\end{block}

\begin{block}{Recursos Online}
\begin{itemize}
    \item \textbf{NCBI:} \url{https://www.ncbi.nlm.nih.gov}
    \item \textbf{PDB:} \url{https://www.rcsb.org}
    \item \textbf{KEGG:} \url{https://www.genome.jp/kegg}
    \item \textbf{BioCyc:} \url{https://biocyc.org}
\end{itemize}
\end{block}
\end{frame}

% ============================================================================
% SLIDE FINAL
% ============================================================================
\begin{frame}
\begin{center}
{\Huge Obrigado!}

\vspace{1cm}

{\Large Capítulo 1: Sistemas Biológicos}

\vspace{1cm}

\textbf{Próximo Capítulo:}\\
Introdução à Modelagem Matemática

\vspace{1cm}

\textit{Dúvidas? Perguntas?}
\end{center}
\end{frame}

\end{document}
